\subsection{章末確認問題}
\begin{itemize}[leftmargin=2em]
\item 1. ソフトウェアセキュリティ、情報セキュリティ、サイバーセキュリティの違いを説明せよ。
\item 2. 機密性、完全性、可用性を、それぞれ一つの失敗例とともに説明せよ。
\item 3. 脆弱性が、ソフトウェア本体だけでなく第三者部品からも入る理由を説明せよ。
\item 4. ISMS がソフトウェアセキュリティ要求を生むことがある理由を説明せよ。
\item 5. アジャイル開発でセキュリティが「ブロッカー」ではなく「実現要因」であるべき理由を説明せよ。
\item 6. SDLC と DevSecOps の違いを説明せよ。
\item 7. セキュリティ要求、セキュリティ設計、セキュリティテストの関係を説明せよ。
\item 8. 静的分析と動的テストの違いを説明せよ。
\item 9. CVE、CWE、CAPEC、CVSS の役割を説明せよ。
\item 10. クラウド、IoT、機械学習アプリケーションのうち一つを選び、その領域固有のセキュリティ上の注意点を説明せよ。
\end{itemize}
\begin{notebox}{解答の方向性}1 は守る対象と範囲の違いを述べる。2 は CIA の三要素を使う。3 はライブラリ、COTS、OS の脆弱性を含める。4 は法令、規制、リスク評価、監査から要求が生じることを述べる。5 は開発速度に合わせた支援、自動化、開発者の責任を述べる。6 はライフサイクル段階への組込みと、開発・運用・セキュリティの継続的統合の違いを述べる。7 以降は、本文の「要求から設計・実装・テスト・脆弱性管理へつなぐ」という観点で説明できる。\end{notebox}
